\subsection{Complete recognition and proof-replay measurements}

The isolated comparison replaces only the geometry validator, loading the
old implementation before timing. Both arms use the same selection wrapper.
Each timed call runs complete recognition with earlier shortcuts disabled,
the native stage enabled, a twelve-second allowance and a 20,000-object
fallback allowance, followed by an independent replay and hash of any native
proof. Native misses continue through the reduced-Khovanov fallback. Cases
marked ``plain'' disable shellings; cases marked ``shell'' enable them. The
old/new comparison within each row keeps that choice fixed.

Five randomized rounds of old, old A/A, new and new A/A calls cover eleven
pipelines after one warm-up round each: 220 measured calls and 44 warm-ups
in 15.901 seconds. A new-arm A/A ratio of 1.053 on the trefoil motivates a
nine-round repeat. It completes 396 measured calls and 44 warm-ups in
29.043 seconds. All 616 measured calls and 88 warm-ups complete correctly.
Within each case, verdict, method, native work count and certificate hash
are identical across all four arms and every round. No tests, audits, builds
or commits overlap either timing run.
\begin{center}
\small
\begin{tabular}{lrrrrr}
\toprule
Pipeline & Old ms & New ms & Old/new & Old A/A & New A/A \\
\midrule
Circle 9, plain & 123.147 & 93.930 & 1.313 & 1.007 & 1.000 \\
Circle 9, shell & 109.032 & 75.259 & 1.425 & 1.007 & 1.004 \\
Random 29, plain & 31.734 & 23.707 & 1.338 & 0.990 & 0.997 \\
Random 29, shell & 57.294 & 41.883 & 1.389 & 1.020 & 0.986 \\
Original disc, plain & 72.279 & 61.827 & 1.158 & 1.002 & 0.999 \\
Original disc, shell & 69.935 & 50.392 & 1.390 & 1.009 & 1.002 \\
Random 142, shell & 141.855 & 114.957 & 1.238 & 0.993 & 0.995 \\
Random 1, shell & 44.022 & 32.310 & 1.357 & 1.002 & 0.988 \\
Random 81, shell & 68.527 & 49.577 & 1.385 & 1.001 & 1.008 \\
trefoil, shell & 24.406 & 20.977 & 1.164 & 0.999 & 0.995 \\
Three-crossing, shell & 39.935 & 28.772 & 1.393 & 1.004 & 1.007 \\
\bottomrule
\end{tabular}

\end{center}
All eleven final paired old/new medians favor the new census, from 1.158 to
1.425. Final A/A medians range from 0.986 to 1.020. Plain raw witnesses improve
by factors of 1.313 and 1.338; the plain optimized-disc fixture improves by
1.158. Shelling pipelines improve by 1.238--1.425 on the positive fixtures,
and the trefoil's complete negative pipeline improves by 1.164. The benefit
comes from repeated geometry construction and verification, while unchanged
source authentication, cocycle extraction and optimization limit the total
gain. Guard counts stay unchanged because guards do not count every union or
array operation.

Absolute timing can drift between cases. In particular, the plain circle
row was measured during a slower interval than the shelling-circle row.
These rows must not be used as an off/on comparison. A separate randomized
experiment measures that choice under the new census, using the preceding
shelling driver unchanged. Its 200 measured calls and 40 warm-ups all finish
correctly in 10.602 seconds, with A/A medians from 0.989 to 1.048. Nine of ten
pipelines still favor shellings off, with paired off/on ratios from 0.577 to
0.857. The original optimized-disc fixture now favors shellings by a factor
of 1.245, but this one favorable case does not justify a general default.
All four added native corpus positives still take longer with shellings than
with their earlier Khovanov fallback in this full-verification experiment.
The shelling option therefore remains disabled by default.

The incidence census is enabled for all finite-manifold geometry preparation;
there is no new recognition option or certificate schema. Its public
validation contract is unchanged, and the ordinary pipeline corpus reaches
no native stage, so these measurements do not establish a default-portfolio
speedup on that corpus. They establish a reduction in geometry work and
measured gains in complete native pipelines, including the independent proof
cost that defeated the previous simplification experiment.

Implementation revision \code{16e918cc7} follows a fresh merge check against
main. Reproduce the validation with the driver
\code{normal\_orbit\_research.link\_census} in \code{audit} mode, and the
paired timing comparison in \code{recognize --rounds 9} mode. The separate
shelling-choice run uses \code{normal\_orbit\_research.shellings} in
\code{recognize --rounds 5} mode. All measured source proofs and discovery
families are unchanged. In particular, this optimization does not supply the
missing general subexponential or quasipolynomial discovery bound. The incoming
unknot inventory and \code{openai/math} revision \code{fd4aeeb2e} remain unchanged.
